\documentclass[12pt,a4paper]{ctexart}

% === 页面设置 ===
\usepackage[top=2.54cm,bottom=2.54cm,left=3.17cm,right=3.17cm]{geometry}
\usepackage{setspace}
\onehalfspacing  % 1.5 倍行距

% === 字体设置 ===
% ctex 默认使用 Windows 字体（宋体、黑体、楷体、仿宋）
\ctexset{
    section/format = \centering\zihao{3}\heiti,          % 章标题：黑体三号居中
    subsection/format = \zihao{4}\heiti,                  % 节标题：黑体四号
    subsubsection/format = \zihao{-4}\heiti,              % 小节标题：黑体小四
}

% === 常用宏包 ===
\usepackage{amsmath,amssymb,amsthm}
\usepackage{graphicx}
\usepackage{booktabs}
\usepackage{multirow}
\usepackage{float}
\usepackage{caption}
\usepackage{subcaption}
\captionsetup{labelsep=space}  % "图 1 xxx" 而非 "图 1: xxx"
\usepackage[hidelinks]{hyperref}
\usepackage{cleveref}
\usepackage{algorithm2e}
\usepackage{enumitem}
\usepackage{tabularx}
\usepackage{longtable}

% === TikZ（仅用于架构图/流程图）===
\usepackage{tikz}
\usetikzlibrary{arrows.meta,positioning,shapes.geometric,fit,calc}
% 注意：pgfplots 仅作为备用，数据图表应使用 \includegraphics 引用 matplotlib 生成的 PDF
\usepackage{pgfplots}
\pgfplotsset{compat=1.18}

% === 引用格式 GB/T 7714 ===
\usepackage[super,sort&compress]{gbt7714}

% === 页眉页脚 ===
\usepackage{fancyhdr}
\pagestyle{fancy}
\fancyhf{}
\fancyhead[C]{\zihao{5}\songti [学校名称]本科毕业论文}
\fancyfoot[C]{\thepage}

% === 定理环境 ===
\newtheorem{definition}{定义}[section]
\newtheorem{theorem}{定理}[section]
\newtheorem{lemma}{引理}[section]
\newtheorem{proposition}{命题}[section]
\newtheorem{corollary}{推论}[section]

\begin{document}

% === 封面 ===
\begin{titlepage}
\centering
\vspace*{2cm}
{\zihao{1}\heiti [学校名称]\par}
\vspace{1cm}
{\zihao{1}\heiti 本科毕业论文\par}
\vspace{3cm}
{\zihao{2}\heiti [论文标题]\par}
\vspace{4cm}
\begin{tabular}{rl}
\zihao{4}\songti 学生姓名：& [姓名] \\
\zihao{4}\songti 学\hspace{2em}号：& [学号] \\
\zihao{4}\songti 指导教师：& [导师姓名] \\
\zihao{4}\songti 学\hspace{2em}院：& [学院名称] \\
\zihao{4}\songti 专\hspace{2em}业：& [专业名称] \\
\end{tabular}
\vfill
{\zihao{4}\songti \today\par}
\end{titlepage}

% === 摘要 ===
\newpage
\begin{center}
{\zihao{3}\heiti 摘\hspace{2em}要}
\end{center}
\vspace{1em}
\zihao{-4}\songti
[中文摘要内容，300-500字]

\vspace{1em}
\noindent\textbf{关键词：}[关键词1]；[关键词2]；[关键词3]

% === 英文摘要 ===
\newpage
\begin{center}
{\zihao{3}\bfseries Abstract}
\end{center}
\vspace{1em}
[English abstract, 200-400 words]

\vspace{1em}
\noindent\textbf{Keywords:} keyword1; keyword2; keyword3

% === 目录 ===
\newpage
\tableofcontents
\newpage

% === 正文 ===
\input{sections/1_introduction}
\input{sections/2_related_work}
\input{sections/3_method}
\input{sections/4_experiments}
\input{sections/5_conclusion}

% === 参考文献 ===
\newpage
\bibliography{references}

% === 致谢 ===
\newpage
\section*{致\hspace{2em}谢}
\addcontentsline{toc}{section}{致谢}
[致谢内容]

% === 附录 ===
\newpage
\appendix
\input{sections/A_appendix}

\end{document}
